\documentclass[10pt,a4paper]{amsart}
\usepackage{amsmath,amssymb,amsthm,graphicx}
% Retained minimal preamble for the Stacks Project chapter varieties.tex.
% Verbatim excerpt of the chapter preamble: Xy-pic setup (native preamble lines 14-18), the native theorem
% declarations with the mathematical macros (native preamble lines 154-217), and the project own no-op
% definitions of the reference/slogan/history display environments (native wrapper.tex lines 1-3) that the
% counted statement uses for its slogan. Copied mechanically; nothing added, removed or redefined.
% For commutative diagrams we use Xy-pic
\usepackage[all]{xy}

% We use 2cell for 2-commutative diagrams.
\xyoption{2cell}
% Theorem environments.
%
\theoremstyle{plain}
\newtheorem{theorem}[subsection]{Theorem}
\newtheorem{proposition}[subsection]{Proposition}
\newtheorem{lemma}[subsection]{Lemma}

\theoremstyle{definition}
\newtheorem{definition}[subsection]{Definition}
\newtheorem{example}[subsection]{Example}
\newtheorem{exercise}[subsection]{Exercise}
\newtheorem{situation}[subsection]{Situation}

\theoremstyle{remark}
\newtheorem{remark}[subsection]{Remark}
\newtheorem{remarks}[subsection]{Remarks}

\numberwithin{equation}{subsection}

% Macros
%
\def\lim{\mathop{\mathrm{lim}}\nolimits}
\def\colim{\mathop{\mathrm{colim}}\nolimits}
\def\Spec{\mathop{\mathrm{Spec}}}
\def\Hom{\mathop{\mathrm{Hom}}\nolimits}
\def\Ext{\mathop{\mathrm{Ext}}\nolimits}
\def\SheafHom{\mathop{\mathcal{H}\!\mathit{om}}\nolimits}
\def\SheafExt{\mathop{\mathcal{E}\!\mathit{xt}}\nolimits}
\def\Sch{\mathit{Sch}}
\def\Mor{\mathop{\mathrm{Mor}}\nolimits}
\def\Ob{\mathop{\mathrm{Ob}}\nolimits}
\def\Sh{\mathop{\mathit{Sh}}\nolimits}
\def\NL{\mathop{N\!L}\nolimits}
\def\CH{\mathop{\mathrm{CH}}\nolimits}
\def\proetale{{pro\text{-}\acute{e}tale}}
\def\etale{{\acute{e}tale}}
\def\QCoh{\mathit{QCoh}}
\def\Ker{\mathop{\mathrm{Ker}}}
\def\Im{\mathop{\mathrm{Im}}}
\def\Coker{\mathop{\mathrm{Coker}}}
\def\Coim{\mathop{\mathrm{Coim}}}

% Boxtimes
%
\DeclareMathSymbol{\boxtimes}{\mathbin}{AMSa}{"02}

%
% Macros for moduli stacks/spaces
%
\def\QCohstack{\mathcal{QC}\!\mathit{oh}}
\def\Cohstack{\mathcal{C}\!\mathit{oh}}
\def\Spacesstack{\mathcal{S}\!\mathit{paces}}
\def\Quotfunctor{\mathrm{Quot}}
\def\Hilbfunctor{\mathrm{Hilb}}
\def\Curvesstack{\mathcal{C}\!\mathit{urves}}
\def\Polarizedstack{\mathcal{P}\!\mathit{olarized}}
\def\Complexesstack{\mathcal{C}\!\mathit{omplexes}}
% \Pic is the operator that assigns to X its picard group, usage \Pic(X)
% \Picardstack_{X/B} denotes the Picard stack of X over B
% \Picardfunctor_{X/B} denotes the Picard functor of X over B
\def\Pic{\mathop{\mathrm{Pic}}\nolimits}
\def\Picardstack{\mathcal{P}\!\mathit{ic}}
\def\Picardfunctor{\mathrm{Pic}}
\def\Deformationcategory{\mathcal{D}\!\mathit{ef}}
\newenvironment{reference}{}{}
\newenvironment{slogan}{}{}
\newenvironment{history}{}{}
\usepackage{cleveref}
\crefname{originalreference}{original reference}{original references}
\Crefname{originalreference}{Original reference}{Original references}
\makeatletter
\newlabel{cohomology-lemma-h1-invertible}{{09NU}{source}{}{sourcepointer-1}{}}
\newlabel{cohomology-lemma-h1-invertible@cref}{{[originalreference][0][]09NU}{[1][0][]source}{}{}{}}
\makeatother
\crefname{originalreference}{original reference}{original references}
\Crefname{originalreference}{Original reference}{Original references}
\makeatletter
\newlabel{cohomology-lemma-Leray}{{01F2}{source}{}{sourcepointer-2}{}}
\newlabel{cohomology-lemma-Leray@cref}{{[originalreference][0][]01F2}{[1][0][]source}{}{}{}}
\makeatother
\crefname{originalreference}{original reference}{original references}
\Crefname{originalreference}{Original reference}{Original references}
\makeatletter
\newlabel{divisors-lemma-finite-trivialize-invertible-upstairs}{{0BUT}{source}{}{sourcepointer-3}{}}
\newlabel{divisors-lemma-finite-trivialize-invertible-upstairs@cref}{{[originalreference][0][]0BUT}{[1][0][]source}{}{}{}}
\makeatother
\crefname{originalreference}{original reference}{original references}
\Crefname{originalreference}{Original reference}{Original references}
\makeatletter
\newlabel{cohomology-lemma-cohomology-and-closed-immersions}{{02UV}{source}{}{sourcepointer-4}{}}
\newlabel{cohomology-lemma-cohomology-and-closed-immersions@cref}{{[originalreference][0][]02UV}{[1][0][]source}{}{}{}}
\makeatother
\crefname{originalreference}{original reference}{original references}
\Crefname{originalreference}{Original reference}{Original references}
\makeatletter
\newlabel{modules-lemma-i-star-exact}{{01AX}{source}{}{sourcepointer-5}{}}
\newlabel{modules-lemma-i-star-exact@cref}{{[originalreference][0][]01AX}{[1][0][]source}{}{}{}}
\makeatother
\title{Every invertible sheaf on the source of a finite morphism that is an isomorphism outside a discrete set is pulled back from the target}
\author{The Stacks Project Authors}
\date{stacks-project:varieties.tex:local-snapshot-20261011; GNU FDL 1.2 or later}
\begin{document}
\maketitle
\noindent\textit{Editorial note.} All delivered mathematical text is the author's own, copied verbatim from the pinned Stacks Project chapter file varieties.tex; no statement or proof step was rewritten, reordered or removed.\par\medskip
\noindent\textit{Editorial note.} The counted span cites five cross-chapter labels, each recorded in assembly.external\_references with its canonical tag as a navigation pointer.\par\medskip
\noindent\textit{Editorial note.} The proof cites no bibliography item inside the counted span.\par\medskip
\noindent\textit{Allowed standard prerequisites (editorial).}\begin{itemize}
\item Picard group as H\textasciicircum{}1 of the unit sheaf and the Leray spectral sequence. Stacks Project, Cohomology, Lemmas cohomology-lemma-h1-invertible (tag 09NU) and cohomology-lemma-Leray (tag 01F2). Convert the Picard statement into a statement about H\textasciicircum{}1 and provide the spectral sequence comparison.
\item Trivializing invertible sheaves after a finite morphism and cohomology supported on a closed subset. Stacks Project, Divisors, Lemma divisors-lemma-finite-trivialize-invertible-upstairs (tag 0BUT); Cohomology, Lemma cohomology-lemma-cohomology-and-closed-immersions (tag 02UV). The first kills the edge map and the second handles sheaves supported on the discrete complement.
\end{itemize}\medskip
\ifdefined\hypertarget\hypertarget{sourcepointer-1}{}\fi
\noindent\textit{Source reference.} Original reference 09NU: sources/stacks-chow/tags.txt, source lines 12227--12227. Cross-chapter standard prerequisite cited by the author inside the counted span; it is used as a pointer and its own proof is not part of this item.\par\medskip
\ifdefined\hypertarget\hypertarget{sourcepointer-2}{}\fi
\noindent\textit{Source reference.} Original reference 01F2: sources/stacks-chow/tags.txt, source lines 1855--1855. Cross-chapter standard prerequisite cited by the author inside the counted span; it is used as a pointer and its own proof is not part of this item.\par\medskip
\ifdefined\hypertarget\hypertarget{sourcepointer-3}{}\fi
\noindent\textit{Source reference.} Original reference 0BUT: sources/stacks-chow/tags.txt, source lines 14886--14886. Cross-chapter standard prerequisite cited by the author inside the counted span; it is used as a pointer and its own proof is not part of this item.\par\medskip
\ifdefined\hypertarget\hypertarget{sourcepointer-4}{}\fi
\noindent\textit{Source reference.} Original reference 02UV: sources/stacks-chow/tags.txt, source lines 3746--3746. Cross-chapter standard prerequisite cited by the author inside the counted span; it is used as a pointer and its own proof is not part of this item.\par\medskip
\ifdefined\hypertarget\hypertarget{sourcepointer-5}{}\fi
\noindent\textit{Source reference.} Original reference 01AX: sources/stacks-chow/tags.txt, source lines 1706--1706. Cross-chapter standard prerequisite cited by the author inside the counted span; it is used as a pointer and its own proof is not part of this item.\par\medskip
\begin{lemma}
\label{lemma-surjective-pic-birational-finite}
Let $f : X \to Y$ be a finite morphism of schemes. Assume there
exists an open $V \subset Y$ such that $f^{-1}(V) \to V$ is an
isomorphism and $Y \setminus V$ is a discrete space. Then every
invertible $\mathcal{O}_X$-module is the pullback of an invertible
$\mathcal{O}_Y$-module.
\end{lemma}

\begin{proof}
We will use that $\Pic(X) = H^1(X, \mathcal{O}_X^*)$, see
Cohomology, Lemma \ref{cohomology-lemma-h1-invertible}.
Consider the Leray spectral sequence for the abelian sheaf $\mathcal{O}_X^*$
and $f$, see Cohomology, Lemma \ref{cohomology-lemma-Leray}.
Consider the induced map
$$
H^1(X, \mathcal{O}_X^*) \longrightarrow H^0(Y, R^1f_*\mathcal{O}_X^*)
$$
Divisors, Lemma \ref{divisors-lemma-finite-trivialize-invertible-upstairs}
says exactly that this map is zero. Hence Leray gives
$H^1(X, \mathcal{O}_X^*) = H^1(Y, f_*\mathcal{O}_X^*)$.
Next we consider the map
$$
f^\sharp : \mathcal{O}_Y^* \longrightarrow f_*\mathcal{O}_X^*
$$
By assumption the kernel and cokernel of this map are supported
on the closed subset $T = Y \setminus V$ of $Y$.
Since $T$ is a discrete topological space by assumption
the higher cohomology groups of any abelian sheaf on $Y$ supported
on $T$ is zero (follows from
Cohomology, Lemma \ref{cohomology-lemma-cohomology-and-closed-immersions},
Modules, Lemma \ref{modules-lemma-i-star-exact}, and
the fact that $H^i(T, \mathcal{F}) = 0$ for any $i > 0$
and any abelian sheaf $\mathcal{F}$ on $T$).
Breaking the displayed map into short exact sequences
$$
0 \to \Ker(f^\sharp) \to \mathcal{O}_Y^* \to \Im(f^\sharp) \to 0,\quad
0 \to \Im(f^\sharp) \to f_*\mathcal{O}_X^* \to \Coker(f^\sharp) \to 0
$$
we first conclude that $H^1(Y, \mathcal{O}_Y^*) \to H^1(Y, \Im(f^\sharp))$
is surjective and then that
$H^1(Y, \Im(f^\sharp)) \to H^1(Y, f_*\mathcal{O}_X^*)$ is surjective.
Combining all the above we find that $H^1(Y, \mathcal{O}_Y^*) \to
H^1(X, \mathcal{O}_X^*)$ is surjective as desired.
\end{proof}

\end{document}
